\section{Main results}
\label{sec:main-achievements}

The first result is a common semantic account of the source and target
languages. \zflean{} (\Cref{ch:zflean}) supplies the ZFC universe in which the B and SMT-LIB
formalizations of \Cref{ch:b,ch:smt} are interpreted. Their syntax, typing, and
denotations can therefore be related within one mathematical model. This turns
the correctness of the encoding from an informal comparison of formulas
into a statement about preservation of meaning.

The second result is that this preservation does not require a single fixed
representation for each B value. Type loosening organizes compatible target
representations into a decidable order and bridges functional representations and more general relational ones while preserving semantics. On this
basis, the central theorem for \beer{} relates every successful encoding of a
supported, well-typed, and well-defined term to the representation selected by
the encoder, including the declarations and scoping information generated
during translation. For Boolean terms, the result yields preservation of
truth. \textsc{Zero-Proof} \beer{} deliberately falls outside this guarantee:
its broader coverage is experimental rather than verified.

The third result concerns the interactive route. \barrel{} makes the partiality
of B operators explicit by requiring well-definedness evidence in their Lean
encoding. Such evidence becomes part of the proof state, so an ill-defined
instantiation cannot silently be used in a completed proof. Proof search may
combine domain-specific lemmas, tactics, and manual steps, with every resulting
theorem being checked by Lean's kernel.

Neither the automated nor the interactive route is a complete verification of a complete B toolchain yet,
but each identifies what has been established and what must still be trusted.
